\section{Experimental Setup}

We validate each system component through targeted experiments testing specific mechanistic claims. Experiments use proof-of-concept scale infrastructure with mock/simulated data to isolate mechanism feasibility from deployment complexity.

\subsection{Experimental Questions}

\begin{enumerate}
\item \textbf{Q1 (h-e1):} Can load-time instrumentation track adoption with <10\% overhead and $\geq$95\% capture rate?
\item \textbf{Q2 (h-m1):} Do health metrics achieve $\geq$60\% precision in identifying deprecation candidates?
\item \textbf{Q3 (h-m2):} Does context-aware inference achieve $\geq$70\% accuracy on task-specific successors?
\item \textbf{Q4 (h-m3):} Can dependency graphs achieve $\geq$95\% impact completeness for migration planning?
\item \textbf{Q5 (h-m4):} Does full telemetry stack maintain $\geq$95\% capture with <10\% overhead?
\end{enumerate}

\subsection{Datasets and Baselines}

\textbf{HuggingFace Datasets Hub Infrastructure:}
\begin{itemize}
\item \textbf{h-e1:} 100 benchmark dataset loads (MNIST, CIFAR-10, ImageNet subsets) to measure overhead
\item \textbf{h-m1:} 1000-dataset simulation with synthetic health metrics (velocity, emergence, issue\_ratio)
\item \textbf{h-m2:} 50 synthetic validation cases with ground-truth task labels
\item \textbf{h-m3:} Mock 4-node dependency graph (dataset $\rightarrow$ downstream models)
\item \textbf{h-m4:} Integrated stack with 500 simulated user sessions
\end{itemize}

\textbf{Baselines:} No formal baselines exist (no prior dataset deprecation systems). Performance targets derived from software package manager requirements (NPM overhead benchmarks, PyPI deprecation warning thresholds).

\textbf{Mock Data Rationale:} Real HuggingFace API integration blocked by configuration errors and rate limits during PoC validation. Mock data isolates mechanism validation from external dependencies; generalization to real API data is future work.

\subsection{Evaluation Metrics}

\begin{itemize}
\item \textbf{Overhead:} $(t_\text{instrumented} - t_\text{baseline}) / t_\text{baseline} \times 100\%$
\item \textbf{Capture Rate:} $\text{captured\_events} / \text{total\_events}$
\item \textbf{Precision:} $\text{TP} / (\text{TP} + \text{FP})$
\item \textbf{Recall:} $\text{TP} / (\text{TP} + \text{FN})$
\item \textbf{Context Accuracy:} $\text{correct\_inferences} / \text{total\_inferences}$
\item \textbf{Override Rate:} $\text{user\_overrides} / \text{total\_recommendations}$
\end{itemize}

Wilson confidence intervals (95\% CI) computed for capture rate to validate statistical reliability.

\subsection{Experimental Protocol}

Each sub-hypothesis follows gate protocol:
\begin{enumerate}
\item \textbf{Setup:} Load mock/simulated data matching experimental question
\item \textbf{Execution:} Run mechanism under test conditions
\item \textbf{Measurement:} Compute target metrics (overhead, precision, accuracy, etc.)
\item \textbf{Gate Check:} Pass if all metrics meet thresholds, otherwise FAIL
\end{enumerate}

MUST\_WORK gates (h-e1, h-m1, h-m2) are critical for system viability. SHOULD\_WORK gates (h-m3, h-m4) validate supporting infrastructure but don't block Phase 6 progression.

\subsection{Reproducibility}

Code, mock data generators, and evaluation scripts available at [repository TBD]. Simulated HuggingFace metadata uses patterns from literature (software package deprecation velocity distributions, GitHub issue accumulation rates). PoC scale chosen to validate mechanism feasibility; production validation requires deployment testing.
